\documentclass[11pt]{article}
\usepackage[margin=1in]{geometry}
\usepackage{amsmath,amssymb}
\usepackage{graphicx}
\usepackage{booktabs}
\usepackage{longtable}
\usepackage{hyperref}
\usepackage{placeins}

\renewcommand{\thesection}{S\arabic{section}}
\renewcommand{\thetable}{S\arabic{table}}
\renewcommand{\thefigure}{S\arabic{figure}}

\title{Supplementary Material\\[4pt]\large Physics-Informed Neural Operator Learning for Hyperelastic Solids: From 2D Benchmarks to a 3D Rocking Rubber-Mount Bushing}
\author{Omar Amro, Timon Rabczuk}
\date{}

\begin{document}
\maketitle

\noindent This supplement reports numerical evidence that supports claims
made in the main paper but was condensed there to a sentence, a figure,
and/or a one-row summary. All values below are reproduced directly from
the project's own verified report
(\texttt{PFEM\_Transolver\_Report\_2026-10-04.docx}); none are estimated,
interpolated, or reconstructed. Section/table numbers here are
self-contained (S1, S2, \ldots) and distinct from the main paper's own
numbering.

\section{Full matched-batch GPU-native-FEM-vs-operator comparison}
\label{sec:s1}
The main paper (Section 6.3) reports only the aggregate speed-up range
($1{,}215$--$1{,}297\times$ at matched batch size) and the resulting
break-even range ($1{,}133$--$95{,}038$ instances) against the GPU-native
FEM baseline. The report explicitly cautions that break-even without a
batch size is misleading, since both the per-sample FEM cost and the
operator's own latency vary substantially with batch size. Tables~\ref{tab:s1-fem}--\ref{tab:s1-breakeven}
report the full per-case, per-batch-size breakdown.

\begin{table}[h]
\centering
\caption{GPU-native batched FEM timing, $N{=}21$, all six cases,
  per-sample cost (ms, median of 3 repeats) at four batch sizes, and
  speed-up at $\mathrm{bs}{=}128$ relative to each case's own directly
  measured CPU reference.}
\label{tab:s1-fem}
\small
\begin{tabular}{lccccc}
\toprule
Case & bs$=$1 & bs$=$8 & bs$=$32 & bs$=$128 & Speed-up @ bs$=$128 \\
\midrule
B1 $\times$ NH & 1{,}651.6 & 477.9 & 381.3 & 354.6 & $71.72\times$ \\
B1 $\times$ MR & 2{,}047.7 & 492.4 & 388.3 & 359.7 & $149.4\times$ \\
B1 $\times$ AB & 2{,}524.6 & 525.5 & 410.9 & 378.3 & $138.9\times$ \\
B2 $\times$ NH & 1{,}665.9 & 479.0 & 381.3 & 354.7 & $73.05\times$ \\
B2 $\times$ MR & 2{,}010.5 & 492.8 & 389.1 & 359.8 & $171.5\times$ \\
B2 $\times$ AB & 2{,}612.1 & 529.3 & 411.1 & 378.5 & $159.3\times$ \\
\bottomrule
\end{tabular}
\end{table}

\begin{table}[h]
\centering
\caption{Trained-operator inference latency (ms/sample), same batch
  sizes as Table~\ref{tab:s1-fem}, all six cases (median of 50 timed
  repeats after 10 warm-up calls, same GPU).}
\label{tab:s1-op}
\small
\begin{tabular}{lcccc}
\toprule
Case & bs$=$1 & bs$=$8 & bs$=$32 & bs$=$128 \\
\midrule
B1 $\times$ NH & 4.582 & 0.573 & 0.336 & 0.292 \\
B1 $\times$ MR & 4.832 & 0.607 & 0.336 & 0.292 \\
B1 $\times$ AB & 4.753 & 0.599 & 0.336 & 0.292 \\
B2 $\times$ NH & 4.680 & 0.585 & 0.336 & 0.291 \\
B2 $\times$ MR & 4.714 & 0.594 & 0.336 & 0.291 \\
B2 $\times$ AB & 4.799 & 0.602 & 0.336 & 0.292 \\
\bottomrule
\end{tabular}
\end{table}

\begin{table}[h]
\centering
\caption{Speed-up of the trained operator over the GPU-native FEM
  solver, both measured at the same batch size (Table~\ref{tab:s1-fem}
  divided by Table~\ref{tab:s1-op} row by row) -- a genuinely
  hardware-matched comparison.}
\label{tab:s1-speedup}
\small
\begin{tabular}{lcccc}
\toprule
Case & bs$=$1 & bs$=$8 & bs$=$32 & bs$=$128 \\
\midrule
B1 $\times$ NH & $360\times$ & $833\times$   & $1{,}134\times$ & $1{,}215\times$ \\
B1 $\times$ MR & $424\times$ & $812\times$   & $1{,}154\times$ & $1{,}232\times$ \\
B1 $\times$ AB & $531\times$ & $878\times$   & $1{,}221\times$ & $1{,}297\times$ \\
B2 $\times$ NH & $356\times$ & $819\times$   & $1{,}135\times$ & $1{,}217\times$ \\
B2 $\times$ MR & $426\times$ & $830\times$   & $1{,}158\times$ & $1{,}235\times$ \\
B2 $\times$ AB & $544\times$ & $879\times$   & $1{,}223\times$ & $1{,}297\times$ \\
\bottomrule
\end{tabular}
\end{table}

\begin{table}[h]
\centering
\caption{Break-even, in new problem instances, against both baselines:
  the CPU-FEM reference (batch size 1 only, since that reference is not
  batched) and the GPU-native FEM solver at each matched batch size.}
\label{tab:s1-breakeven}
\small
\begin{tabular}{lccccc}
\toprule
Case & vs.\ CPU FEM & \multicolumn{4}{c}{vs.\ GPU-native FEM} \\
 & (bs$=$1) & bs$=$1 & bs$=$8 & bs$=$32 & bs$=$128 \\
\midrule
B1 $\times$ NH & 113   & 1{,}745  & 6{,}021  & 7{,}543  & 8{,}112 \\
B1 $\times$ MR & 52    & 1{,}363  & 5{,}663  & 7{,}179  & 7{,}751 \\
B1 $\times$ AB & 54    & 1{,}133  & 5{,}441  & 6{,}956  & 7{,}554 \\
B2 $\times$ NH & 1{,}245 & 19{,}410 & 67{,}391 & 84{,}627 & 90{,}990 \\
B2 $\times$ MR & 554   & 17{,}033 & 69{,}404 & 87{,}884 & 95{,}038 \\
B2 $\times$ AB & 412   & 9{,}530  & 46{,}993 & 60{,}490 & 65{,}698 \\
\bottomrule
\end{tabular}
\end{table}

The GPU-break-even figures span $1{,}133$ (B1$\times$AB, bs$=$1) to
$95{,}038$ (B2$\times$MR, bs$=$128) -- the full range the main paper
quotes only as an aggregate; which end of that range applies depends
entirely on the batch size assumed practical for the deployment
scenario, exactly as the main paper's own cautionary sentence states.

\FloatBarrier
\section{Full 16-resolution six-case FEM-vs-operator QoI sweeps}
\label{sec:s2}
Section 6.6 of the main paper reports only the coarsest finite-element
mesh reaching a fixed $5\%$ threshold per quantity of interest (Table 20)
and the binding-quantity accuracy-matched cost comparison (Table 21).
The report's own per-case sweeps cover sixteen resolutions directly
($N{=}3$, $4$, $5$, $6$, $9$, $11$, $13$, $17$, $21$, $25$, $29$, $33$,
$37$, $41$, $45$, $49$), FEM and operator both
solved on the identical field realization and scored against the same
$N{=}201$ fine reference; the full sweeps are reported here, two tables
per case (displacement/$H^1$/energy/reaction, and fixed-region Cauchy
stress average/99th-percentile/true-max). B2's reaction cells are n/a
for both methods in every row, exactly as in the source (no single
reaction-resultant metric is defined for the B2 ring geometry, main
paper Table 20's caption).

\begin{table}[h]
\centering
\caption{B1$\times$Neo-Hookean, FEM vs.\ operator, all sixteen resolutions:
  displacement $L^2$, $H^1$ semi-norm, tangent energy, reaction.}
\label{tab:s2-b1nh-a}
\scriptsize
\begin{tabular}{lcccccccc}
\toprule
$N$ & FEM $L^2$ & Op.\ $L^2$ & FEM $H^1$ & Op.\ $H^1$ & FEM En.\ & Op.\ En.\ & FEM Rct.\ & Op.\ Rct.\ \\
\midrule
3  & 6.78\%  & 31.06\% & 21.79\% & 42.62\% & 18.63\% & 35.19\% & 12.00\% & 39.24\% \\
4  & 3.63\%  & 17.00\% & 16.06\% & 32.71\% & 13.53\% & 28.30\% & 5.94\%  & 18.98\% \\
5  & 2.31\%  & 11.11\% & 12.91\% & 30.06\% & 10.88\% & 26.51\% & 3.39\%  & 4.12\%  \\
6  & 1.64\%  & 10.46\% & 10.67\% & 28.18\% & 9.17\%  & 23.49\% & 2.17\%  & 4.81\%  \\
9  & 0.81\%  & 8.81\%  & 7.14\%  & 24.27\% & 6.43\%  & 14.96\% & 0.84\%  & 10.13\% \\
11 & 0.58\%  & 6.98\%  & 5.90\%  & 22.36\% & 5.45\%  & 12.50\% & 0.54\%  & 8.52\%  \\
13 & 0.45\%  & 5.98\%  & 5.43\%  & 20.91\% & 4.63\%  & 11.12\% & 0.38\%  & 7.01\%  \\
17 & 0.29\%  & 4.40\%  & 4.36\%  & 17.80\% & 3.72\%  & 9.10\%  & 0.22\%  & 4.84\%  \\
21 & 0.21\%  & 3.43\%  & 3.19\%  & 15.46\% & 3.25\%  & 7.88\%  & 0.14\%  & 3.59\%  \\
25 & 0.16\%  & 2.78\%  & 3.36\%  & 13.90\% & 2.69\%  & 6.90\%  & 0.10\%  & 2.82\%  \\
29 & 0.13\%  & 2.33\%  & 2.84\%  & 12.52\% & 2.39\%  & 6.26\%  & 0.07\%  & 2.29\%  \\
33 & 0.10\%  & 2.03\%  & 2.38\%  & 11.40\% & 2.13\%  & 5.75\%  & 0.06\%  & 1.97\%  \\
37 & 0.09\%  & 1.83\%  & 2.16\%  & 10.52\% & 1.91\%  & 5.35\%  & 0.04\%  & 1.77\%  \\
41 & 0.07\%  & 1.72\%  & 1.56\%  & 9.71\%  & 1.90\%  & 5.12\%  & 0.04\%  & 1.64\%  \\
45 & 0.06\%  & 1.68\%  & 2.66\%  & 9.39\%  & 1.60\%  & 4.78\%  & 0.03\%  & 1.56\%  \\
49 & 0.06\%  & 1.67\%  & 2.35\%  & 8.82\%  & 1.50\%  & 4.58\%  & 0.02\%  & 1.51\%  \\
\bottomrule
\end{tabular}
\end{table}

\begin{table}[h]
\centering
\caption{B1$\times$Neo-Hookean, fixed-region Cauchy-stress error
  (region-weighted average, 99th percentile, true max), FEM vs.\
  operator, all sixteen resolutions.}
\label{tab:s2-b1nh-b}
\scriptsize
\begin{tabular}{lcccccc}
\toprule
$N$ & FEM avg & Op.\ avg & FEM p99 & Op.\ p99 & FEM max & Op.\ max \\
\midrule
3  & 3.86\% & 19.58\% & 27.74\% & 43.96\% & 64.55\% & 72.51\% \\
4  & 5.91\% & 5.35\%  & 24.98\% & 33.46\% & 63.08\% & 67.31\% \\
5  & 7.07\% & 3.54\%  & 21.84\% & 25.82\% & 61.28\% & 63.37\% \\
6  & 7.46\% & 11.53\% & 18.71\% & 17.71\% & 59.48\% & 59.13\% \\
9  & 3.77\% & 16.21\% & 10.74\% & 2.02\%  & 54.75\% & 50.46\% \\
11 & 1.73\% & 11.70\% & 6.82\%  & 1.95\%  & 52.08\% & 47.81\% \\
13 & 1.57\% & 9.08\%  & 3.13\%  & 4.96\%  & 49.71\% & 45.77\% \\
17 & 1.27\% & 5.56\%  & 2.47\%  & 9.35\%  & 45.62\% & 42.08\% \\
21 & 0.92\% & 3.89\%  & 6.21\%  & 13.26\% & 42.17\% & 38.62\% \\
25 & 0.63\% & 2.81\%  & 9.21\%  & 15.86\% & 39.19\% & 35.84\% \\
29 & 0.61\% & 2.17\%  & 11.45\% & 17.50\% & 36.54\% & 33.40\% \\
33 & 0.37\% & 1.60\%  & 12.74\% & 18.65\% & 34.17\% & 31.21\% \\
37 & 0.44\% & 1.33\%  & 13.36\% & 19.49\% & 32.03\% & 29.24\% \\
41 & 0.35\% & 0.96\%  & 13.51\% & 20.13\% & 30.06\% & 27.40\% \\
45 & 0.35\% & 0.72\%  & 13.34\% & 20.51\% & 28.26\% & 25.66\% \\
49 & 0.30\% & 0.45\%  & 12.02\% & 19.80\% & 26.58\% & 23.96\% \\
\bottomrule
\end{tabular}
\end{table}

\begin{table}[h]
\centering
\caption{B1$\times$Mooney-Rivlin, FEM vs.\ operator, all sixteen
  resolutions: displacement $L^2$, $H^1$ semi-norm, tangent energy,
  reaction.}
\label{tab:s2-b1mr-a}
\scriptsize
\begin{tabular}{lcccccccc}
\toprule
$N$ & FEM $L^2$ & Op.\ $L^2$ & FEM $H^1$ & Op.\ $H^1$ & FEM En.\ & Op.\ En.\ & FEM Rct.\ & Op.\ Rct.\ \\
\midrule
3  & 6.55\% & 37.02\% & 21.20\% & 67.78\% & 18.20\% & 49.16\% & 9.67\% & 18.10\% \\
4  & 3.51\% & 24.51\% & 15.87\% & 49.16\% & 13.43\% & 29.24\% & 4.86\% & 6.21\%  \\
5  & 2.25\% & 19.21\% & 12.86\% & 38.81\% & 10.87\% & 22.56\% & 2.80\% & 3.22\%  \\
6  & 1.61\% & 15.91\% & 10.68\% & 34.45\% & 9.20\%  & 19.62\% & 1.81\% & 1.65\%  \\
9  & 0.80\% & 12.26\% & 7.21\%  & 27.17\% & 6.49\%  & 14.82\% & 0.72\% & 3.11\%  \\
11 & 0.58\% & 11.14\% & 5.97\%  & 24.71\% & 5.51\%  & 12.93\% & 0.46\% & 2.88\%  \\
13 & 0.45\% & 10.39\% & 5.50\%  & 22.85\% & 4.71\%  & 11.68\% & 0.32\% & 2.67\%  \\
17 & 0.30\% & 9.47\%  & 4.43\%  & 20.14\% & 3.79\%  & 10.00\% & 0.19\% & 2.20\%  \\
21 & 0.21\% & 8.88\%  & 3.25\%  & 18.13\% & 3.31\%  & 8.95\%  & 0.12\% & 1.82\%  \\
25 & 0.16\% & 8.47\%  & 3.43\%  & 16.81\% & 2.76\%  & 8.07\%  & 0.08\% & 1.46\%  \\
29 & 0.13\% & 8.15\%  & 2.90\%  & 15.64\% & 2.45\%  & 7.46\%  & 0.06\% & 1.18\%  \\
33 & 0.11\% & 7.90\%  & 2.43\%  & 14.68\% & 2.18\%  & 6.97\%  & 0.05\% & 0.95\%  \\
37 & 0.09\% & 7.69\%  & 2.20\%  & 13.93\% & 1.97\%  & 6.57\%  & 0.04\% & 0.77\%  \\
41 & 0.08\% & 7.51\%  & 1.60\%  & 13.22\% & 1.95\%  & 6.33\%  & 0.03\% & 0.66\%  \\
45 & 0.07\% & 7.35\%  & 2.73\%  & 12.89\% & 1.65\%  & 6.01\%  & 0.03\% & 0.62\%  \\
49 & 0.06\% & 7.21\%  & 2.41\%  & 12.38\% & 1.55\%  & 5.80\%  & 0.02\% & 0.66\%  \\
\bottomrule
\end{tabular}
\end{table}

\begin{table}[h]
\centering
\caption{B1$\times$Mooney-Rivlin, fixed-region Cauchy-stress error, FEM
  vs.\ operator, all sixteen resolutions.}
\label{tab:s2-b1mr-b}
\scriptsize
\begin{tabular}{lcccccc}
\toprule
$N$ & FEM avg & Op.\ avg & FEM p99 & Op.\ p99 & FEM max & Op.\ max \\
\midrule
3  & 8.15\% & 19.32\% & 25.24\% & 17.53\% & 64.01\% & 60.16\% \\
4  & 8.43\% & 29.63\% & 23.60\% & 7.95\%  & 63.13\% & 55.33\% \\
5  & 8.84\% & 21.79\% & 20.75\% & 11.92\% & 61.47\% & 57.15\% \\
6  & 8.73\% & 23.90\% & 17.76\% & 6.99\%  & 59.76\% & 54.51\% \\
9  & 4.19\% & 17.64\% & 9.93\%  & 0.75\%  & 55.14\% & 49.85\% \\
11 & 2.00\% & 13.11\% & 5.96\%  & 5.07\%  & 52.50\% & 46.98\% \\
13 & 1.79\% & 10.69\% & 2.43\%  & 8.76\%  & 50.14\% & 44.46\% \\
17 & 1.42\% & 8.37\%  & 3.30\%  & 14.17\% & 46.05\% & 40.13\% \\
21 & 1.02\% & 6.76\%  & 7.11\%  & 17.53\% & 42.59\% & 36.54\% \\
25 & 0.70\% & 5.84\%  & 9.81\%  & 18.94\% & 39.58\% & 33.92\% \\
29 & 0.67\% & 5.23\%  & 11.92\% & 19.33\% & 36.91\% & 32.03\% \\
33 & 0.40\% & 4.59\%  & 13.42\% & 19.08\% & 34.52\% & 30.59\% \\
37 & 0.49\% & 4.33\%  & 14.15\% & 18.42\% & 32.35\% & 29.45\% \\
41 & 0.38\% & 3.89\%  & 14.19\% & 17.46\% & 30.36\% & 28.53\% \\
45 & 0.39\% & 3.63\%  & 13.79\% & 16.00\% & 28.53\% & 27.79\% \\
49 & 0.33\% & 3.29\%  & 12.39\% & 13.14\% & 26.84\% & 27.18\% \\
\bottomrule
\end{tabular}
\end{table}

\begin{table}[h]
\centering
\caption{B1$\times$Arruda-Boyce, FEM vs.\ operator, all sixteen
  resolutions: displacement $L^2$, $H^1$ semi-norm, tangent energy,
  reaction.}
\label{tab:s2-b1ab-a}
\scriptsize
\begin{tabular}{lcccccccc}
\toprule
$N$ & FEM $L^2$ & Op.\ $L^2$ & FEM $H^1$ & Op.\ $H^1$ & FEM En.\ & Op.\ En.\ & FEM Rct.\ & Op.\ Rct.\ \\
\midrule
3  & 5.90\% & 17.82\% & 20.72\% & 36.73\% & 16.85\% & 30.19\% & 12.41\% & 20.61\% \\
4  & 3.05\% & 12.24\% & 14.91\% & 30.79\% & 11.97\% & 25.52\% & 6.00\%  & 7.02\%  \\
5  & 1.88\% & 10.88\% & 11.77\% & 26.41\% & 9.51\%  & 18.63\% & 3.36\%  & 5.19\%  \\
6  & 1.31\% & 10.52\% & 9.60\%  & 25.21\% & 7.93\%  & 16.67\% & 2.12\%  & 7.05\%  \\
9  & 0.62\% & 10.98\% & 6.27\%  & 24.83\% & 5.45\%  & 14.27\% & 0.81\%  & 8.59\%  \\
11 & 0.44\% & 10.10\% & 5.13\%  & 23.28\% & 4.56\%  & 12.85\% & 0.51\%  & 7.56\%  \\
13 & 0.33\% & 9.22\%  & 4.67\%  & 21.66\% & 3.84\%  & 11.68\% & 0.36\%  & 6.53\%  \\
17 & 0.21\% & 7.71\%  & 3.69\%  & 18.48\% & 3.04\%  & 9.91\%  & 0.20\%  & 5.05\%  \\
21 & 0.15\% & 6.73\%  & 2.68\%  & 15.85\% & 2.64\%  & 8.62\%  & 0.13\%  & 4.10\%  \\
25 & 0.11\% & 6.14\%  & 2.79\%  & 14.09\% & 2.16\%  & 7.58\%  & 0.09\%  & 3.53\%  \\
29 & 0.09\% & 5.76\%  & 2.37\%  & 12.72\% & 1.90\%  & 6.87\%  & 0.07\%  & 3.17\%  \\
33 & 0.07\% & 5.51\%  & 1.98\%  & 11.69\% & 1.68\%  & 6.32\%  & 0.05\%  & 2.93\%  \\
37 & 0.06\% & 5.35\%  & 1.81\%  & 10.96\% & 1.50\%  & 5.91\%  & 0.04\%  & 2.78\%  \\
41 & 0.05\% & 5.24\%  & 1.27\%  & 10.36\% & 1.51\%  & 5.64\%  & 0.03\%  & 2.65\%  \\
45 & 0.04\% & 5.17\%  & 2.14\%  & 10.11\% & 1.25\%  & 5.35\%  & 0.03\%  & 2.55\%  \\
49 & 0.04\% & 5.12\%  & 1.89\%  & 9.74\%  & 1.16\%  & 5.16\%  & 0.02\%  & 2.47\%  \\
\bottomrule
\end{tabular}
\end{table}

\begin{table}[h]
\centering
\caption{B1$\times$Arruda-Boyce, fixed-region Cauchy-stress error, FEM
  vs.\ operator, all sixteen resolutions.}
\label{tab:s2-b1ab-b}
\scriptsize
\begin{tabular}{lcccccc}
\toprule
$N$ & FEM avg & Op.\ avg & FEM p99 & Op.\ p99 & FEM max & Op.\ max \\
\midrule
3  & 1.83\% & 7.83\%  & 29.05\% & 33.46\% & 60.34\% & 62.81\% \\
4  & 0.93\% & 0.45\%  & 26.37\% & 26.95\% & 58.81\% & 59.15\% \\
5  & 2.65\% & 19.75\% & 23.54\% & 12.08\% & 56.98\% & 50.64\% \\
6  & 3.67\% & 20.52\% & 20.72\% & 10.13\% & 55.17\% & 49.41\% \\
9  & 2.24\% & 20.20\% & 13.71\% & 0.10\%  & 50.52\% & 42.76\% \\
11 & 0.86\% & 15.06\% & 9.91\%  & 2.29\%  & 47.96\% & 41.06\% \\
13 & 0.80\% & 12.19\% & 6.83\%  & 3.25\%  & 45.70\% & 39.94\% \\
17 & 0.71\% & 9.68\%  & 1.83\%  & 6.09\%  & 41.85\% & 37.30\% \\
21 & 0.53\% & 8.31\%  & 1.71\%  & 9.13\%  & 38.64\% & 34.36\% \\
25 & 0.35\% & 7.57\%  & 4.29\%  & 11.26\% & 35.87\% & 31.85\% \\
29 & 0.35\% & 7.24\%  & 6.39\%  & 12.97\% & 33.44\% & 29.83\% \\
33 & 0.21\% & 6.89\%  & 7.83\%  & 14.16\% & 31.27\% & 28.09\% \\
37 & 0.24\% & 6.79\%  & 8.93\%  & 15.07\% & 29.31\% & 26.58\% \\
41 & 0.20\% & 6.60\%  & 9.70\%  & 15.73\% & 27.52\% & 25.26\% \\
45 & 0.18\% & 6.49\%  & 9.29\%  & 15.85\% & 25.88\% & 24.12\% \\
49 & 0.17\% & 6.39\%  & 8.52\%  & 15.27\% & 24.36\% & 23.13\% \\
\bottomrule
\end{tabular}
\end{table}

\begin{table}[h]
\centering
\caption{B2$\times$Neo-Hookean, FEM vs.\ operator, all sixteen
  resolutions: displacement $L^2$, $H^1$ semi-norm, tangent energy
  (reaction: n/a, no single resultant defined for the B2 ring
  geometry).}
\label{tab:s2-b2nh-a}
\scriptsize
\begin{tabular}{lcccccc}
\toprule
$N$ & FEM $L^2$ & Op.\ $L^2$ & FEM $H^1$ & Op.\ $H^1$ & FEM En.\ & Op.\ En.\ \\
\midrule
3  & 16.15\% & 210.82\% & 35.96\% & 208.44\% & 46.61\% & 280.89\% \\
4  & 8.76\%  & 192.07\% & 27.08\% & 187.88\% & 36.34\% & 224.90\% \\
5  & 5.45\%  & 137.29\% & 20.56\% & 135.92\% & 25.60\% & 150.27\% \\
6  & 3.69\%  & 128.61\% & 16.30\% & 135.14\% & 22.31\% & 133.15\% \\
9  & 1.54\%  & 110.87\% & 9.86\%  & 111.63\% & 9.41\%  & 109.69\% \\
11 & 1.00\%  & 80.01\%  & 7.78\%  & 78.88\%  & 8.55\%  & 80.38\%  \\
13 & 0.70\%  & 52.65\%  & 6.68\%  & 52.38\%  & 5.73\%  & 51.25\%  \\
17 & 0.40\%  & 16.61\%  & 4.89\%  & 19.21\%  & 4.04\%  & 16.67\%  \\
21 & 0.26\%  & 6.29\%   & 3.53\%  & 12.74\%  & 3.55\%  & 9.06\%   \\
25 & 0.18\%  & 14.47\%  & 3.24\%  & 18.96\%  & 2.50\%  & 16.25\%  \\
29 & 0.13\%  & 3.87\%   & 2.86\%  & 10.39\%  & 2.08\%  & 7.10\%   \\
33 & 0.10\%  & 4.40\%   & 2.43\%  & 10.29\%  & 1.77\%  & 6.90\%   \\
37 & 0.08\%  & 7.94\%   & 2.23\%  & 13.24\%  & 1.53\%  & 8.67\%   \\
41 & 0.06\%  & 9.76\%   & 1.37\%  & 15.22\%  & 1.58\%  & 9.87\%   \\
45 & 0.05\%  & 11.52\%  & 1.85\%  & 17.99\%  & 1.19\%  & 8.86\%   \\
49 & 0.04\%  & 18.92\%  & 1.65\%  & 26.05\%  & 1.06\%  & 11.92\%  \\
\bottomrule
\end{tabular}
\end{table}

\begin{table}[h]
\centering
\caption{B2$\times$Neo-Hookean, fixed-region Cauchy-stress error, FEM
  vs.\ operator, all sixteen resolutions.}
\label{tab:s2-b2nh-b}
\scriptsize
\begin{tabular}{lcccccc}
\toprule
$N$ & FEM avg & Op.\ avg & FEM p99 & Op.\ p99 & FEM max & Op.\ max \\
\midrule
3  & 15.29\% & 245.59\% & 14.58\% & 230.26\% & 12.40\% & 234.38\% \\
4  & 10.73\% & 141.72\% & 8.68\%  & 121.41\% & 6.06\%  & 123.02\% \\
5  & 6.61\%  & 101.86\% & 7.71\%  & 101.71\% & 5.81\%  & 106.19\% \\
6  & 4.20\%  & 84.01\%  & 4.22\%  & 85.95\%  & 1.93\%  & 89.38\%  \\
9  & 1.63\%  & 94.24\%  & 1.69\%  & 89.65\%  & 0.07\%  & 91.49\%  \\
11 & 1.17\%  & 67.61\%  & 1.31\%  & 62.71\%  & 0.06\%  & 63.39\%  \\
13 & 0.83\%  & 41.20\%  & 0.97\%  & 38.78\%  & 0.06\%  & 39.69\%  \\
17 & 0.46\%  & 5.65\%   & 0.49\%  & 2.84\%   & 0.03\%  & 3.04\%   \\
21 & 0.27\%  & 10.35\%  & 0.25\%  & 11.97\%  & 0.12\%  & 11.62\%  \\
25 & 0.20\%  & 18.80\%  & 0.17\%  & 19.77\%  & 0.07\%  & 19.60\%  \\
29 & 0.13\%  & 4.24\%   & 0.14\%  & 5.64\%   & 0.11\%  & 3.10\%   \\
33 & 0.11\%  & 4.46\%   & 0.08\%  & 5.52\%   & 0.04\%  & 3.81\%   \\
37 & 0.09\%  & 6.92\%   & 0.06\%  & 8.55\%   & 0.01\%  & 7.16\%   \\
41 & 0.07\%  & 7.88\%   & 0.05\%  & 9.64\%   & 0.06\%  & 8.66\%   \\
45 & 0.06\%  & 4.18\%   & 0.05\%  & 5.94\%   & 0.03\%  & 5.56\%   \\
49 & 0.05\%  & 3.22\%   & 0.04\%  & 2.53\%   & 0.03\%  & 3.23\%   \\
\bottomrule
\end{tabular}
\end{table}

\begin{table}[h]
\centering
\caption{B2$\times$Mooney-Rivlin, FEM vs.\ operator, all sixteen
  resolutions: displacement $L^2$, $H^1$ semi-norm, tangent energy.}
\label{tab:s2-b2mr-a}
\scriptsize
\begin{tabular}{lcccccc}
\toprule
$N$ & FEM $L^2$ & Op.\ $L^2$ & FEM $H^1$ & Op.\ $H^1$ & FEM En.\ & Op.\ En.\ \\
\midrule
3  & 16.02\% & 14.66\% & 35.68\% & 45.72\% & 50.20\% & 65.32\% \\
4  & 8.56\%  & 12.86\% & 26.81\% & 40.61\% & 39.14\% & 53.44\% \\
5  & 5.29\%  & 12.65\% & 20.41\% & 38.67\% & 27.31\% & 41.90\% \\
6  & 3.57\%  & 12.61\% & 16.22\% & 37.29\% & 23.78\% & 39.01\% \\
9  & 1.49\%  & 12.58\% & 9.84\%  & 35.68\% & 9.75\%  & 29.48\% \\
11 & 0.97\%  & 12.64\% & 7.78\%  & 35.24\% & 8.91\%  & 29.05\% \\
13 & 0.68\%  & 12.63\% & 6.67\%  & 34.96\% & 5.90\%  & 27.92\% \\
17 & 0.38\%  & 12.66\% & 4.89\%  & 34.61\% & 4.12\%  & 27.38\% \\
21 & 0.25\%  & 12.68\% & 3.53\%  & 34.41\% & 3.64\%  & 27.27\% \\
25 & 0.17\%  & 12.69\% & 3.24\%  & 34.44\% & 2.54\%  & 27.06\% \\
29 & 0.12\%  & 12.70\% & 2.86\%  & 34.39\% & 2.11\%  & 27.00\% \\
33 & 0.10\%  & 12.71\% & 2.43\%  & 34.38\% & 1.80\%  & 26.97\% \\
37 & 0.07\%  & 12.71\% & 2.23\%  & 34.37\% & 1.55\%  & 26.94\% \\
41 & 0.06\%  & 12.72\% & 1.37\%  & 34.33\% & 1.60\%  & 26.95\% \\
45 & 0.05\%  & 12.72\% & 1.86\%  & 34.35\% & 1.20\%  & 26.92\% \\
49 & 0.04\%  & 12.73\% & 1.65\%  & 34.34\% & 1.08\%  & 26.93\% \\
\bottomrule
\end{tabular}
\end{table}

\begin{table}[h]
\centering
\caption{B2$\times$Mooney-Rivlin, fixed-region Cauchy-stress error, FEM
  vs.\ operator, all sixteen resolutions.}
\label{tab:s2-b2mr-b}
\scriptsize
\begin{tabular}{lcccccc}
\toprule
$N$ & FEM avg & Op.\ avg & FEM p99 & Op.\ p99 & FEM max & Op.\ max \\
\midrule
3  & 12.90\% & 8.56\% & 10.54\% & 6.70\% & 7.75\% & 4.47\% \\
4  & 10.08\% & 4.72\% & 7.62\%  & 5.49\% & 4.88\% & 2.93\% \\
5  & 5.72\%  & 1.85\% & 5.48\%  & 4.32\% & 3.24\% & 3.05\% \\
6  & 3.91\%  & 0.47\% & 3.76\%  & 5.14\% & 1.50\% & 3.80\% \\
9  & 1.51\%  & 3.12\% & 1.05\%  & 2.02\% & 0.39\% & 1.67\% \\
11 & 1.05\%  & 2.41\% & 0.67\%  & 1.43\% & 0.35\% & 1.09\% \\
13 & 0.72\%  & 1.87\% & 0.43\%  & 1.17\% & 0.19\% & 0.76\% \\
17 & 0.42\%  & 1.34\% & 0.10\%  & 0.95\% & 0.27\% & 0.48\% \\
21 & 0.25\%  & 0.95\% & 0.02\%  & 0.92\% & 0.31\% & 0.39\% \\
25 & 0.18\%  & 0.67\% & 0.04\%  & 0.94\% & 0.23\% & 0.37\% \\
29 & 0.12\%  & 0.45\% & 0.07\%  & 0.97\% & 0.22\% & 0.21\% \\
33 & 0.09\%  & 0.37\% & 0.04\%  & 1.00\% & 0.13\% & 0.36\% \\
37 & 0.08\%  & 0.20\% & 0.05\%  & 1.03\% & 0.13\% & 0.46\% \\
41 & 0.06\%  & 0.15\% & 0.02\%  & 1.06\% & 0.15\% & 0.52\% \\
45 & 0.05\%  & 0.04\% & 0.01\%  & 1.09\% & 0.11\% & 0.56\% \\
49 & 0.05\%  & 0.00\% & 0.00\%  & 1.11\% & 0.09\% & 0.59\% \\
\bottomrule
\end{tabular}
\end{table}

\begin{table}[h]
\centering
\caption{B2$\times$Arruda-Boyce, FEM vs.\ operator, all sixteen
  resolutions: displacement $L^2$, $H^1$ semi-norm, tangent energy.}
\label{tab:s2-b2ab-a}
\scriptsize
\begin{tabular}{lcccccc}
\toprule
$N$ & FEM $L^2$ & Op.\ $L^2$ & FEM $H^1$ & Op.\ $H^1$ & FEM En.\ & Op.\ En.\ \\
\midrule
3  & 16.04\% & 133.64\% & 36.30\% & 181.39\% & 39.86\% & 192.10\% \\
4  & 8.84\%  & 115.70\% & 27.49\% & 220.68\% & 31.12\% & 217.84\% \\
5  & 5.56\%  & 62.13\%  & 20.82\% & 84.30\%  & 22.46\% & 66.08\%  \\
6  & 3.78\%  & 76.99\%  & 16.47\% & 91.69\%  & 19.57\% & 78.68\%  \\
9  & 1.59\%  & 30.09\%  & 9.90\%  & 39.20\%  & 8.84\%  & 31.02\%  \\
11 & 1.04\%  & 26.73\%  & 7.80\%  & 33.93\%  & 7.89\%  & 26.27\%  \\
13 & 0.73\%  & 21.56\%  & 6.68\%  & 27.70\%  & 5.47\%  & 20.15\%  \\
17 & 0.41\%  & 7.05\%   & 4.89\%  & 11.51\%  & 3.92\%  & 9.55\%   \\
21 & 0.26\%  & 4.12\%   & 3.52\%  & 8.50\%   & 3.39\%  & 5.94\%   \\
25 & 0.18\%  & 5.22\%   & 3.23\%  & 10.68\%  & 2.44\%  & 6.42\%   \\
29 & 0.13\%  & 4.79\%   & 2.85\%  & 10.38\%  & 2.03\%  & 6.45\%   \\
33 & 0.10\%  & 9.95\%   & 2.42\%  & 11.38\%  & 1.74\%  & 10.05\%  \\
37 & 0.08\%  & 18.78\%  & 2.23\%  & 18.35\%  & 1.50\%  & 18.16\%  \\
41 & 0.06\%  & 25.79\%  & 1.37\%  & 24.90\%  & 1.54\%  & 24.73\%  \\
45 & 0.05\%  & 32.73\%  & 1.85\%  & 31.67\%  & 1.16\%  & 31.32\%  \\
49 & 0.04\%  & 38.86\%  & 1.64\%  & 37.59\%  & 1.04\%  & 37.14\%  \\
\bottomrule
\end{tabular}
\end{table}

\begin{table}[h]
\centering
\caption{B2$\times$Arruda-Boyce, fixed-region Cauchy-stress error, FEM
  vs.\ operator, all sixteen resolutions.}
\label{tab:s2-b2ab-b}
\scriptsize
\begin{tabular}{lcccccc}
\toprule
$N$ & FEM avg & Op.\ avg & FEM p99 & Op.\ p99 & FEM max & Op.\ max \\
\midrule
3  & 18.44\% & 54.10\% & 19.25\% & 191.35\% & 17.64\% & 199.45\% \\
4  & 11.85\% & 44.05\% & 11.02\% & 99.72\%  & 8.66\%  & 106.31\% \\
5  & 7.81\%  & 28.94\% & 10.29\% & 26.39\%  & 8.68\%  & 28.03\%  \\
6  & 4.72\%  & 48.94\% & 5.39\%  & 38.78\%  & 3.21\%  & 40.19\%  \\
9  & 1.81\%  & 16.61\% & 2.64\%  & 10.04\%  & 0.75\%  & 10.84\%  \\
11 & 1.35\%  & 13.16\% & 2.12\%  & 8.14\%   & 0.60\%  & 8.65\%   \\
13 & 0.98\%  & 9.08\%  & 1.74\%  & 5.31\%   & 0.61\%  & 5.69\%   \\
17 & 0.50\%  & 2.40\%  & 1.04\%  & 1.01\%   & 0.42\%  & 1.73\%   \\
21 & 0.31\%  & 4.01\%  & 0.78\%  & 4.14\%   & 0.23\%  & 3.46\%   \\
25 & 0.22\%  & 0.79\%  & 0.52\%  & 0.76\%   & 0.21\%  & 0.30\%   \\
29 & 0.15\%  & 1.69\%  & 0.41\%  & 1.50\%   & 0.11\%  & 1.23\%   \\
33 & 0.13\%  & 11.09\% & 0.36\%  & 10.73\%  & 0.15\%  & 10.43\%  \\
37 & 0.10\%  & 19.02\% & 0.31\%  & 18.61\%  & 0.16\%  & 18.39\%  \\
41 & 0.08\%  & 25.26\% & 0.28\%  & 25.04\%  & 0.10\%  & 24.70\%  \\
45 & 0.08\%  & 31.64\% & 0.22\%  & 31.38\%  & 0.11\%  & 31.05\%  \\
49 & 0.05\%  & 37.25\% & 0.21\%  & 36.86\%  & 0.08\%  & 36.60\%  \\
\bottomrule
\end{tabular}
\end{table}

A pattern visible across all six cases, stated plainly rather than left
for the reader to notice: the operator's error is non-monotone in $N$
for several quantities (most visibly B2's $L^2$/$H^1$/energy columns,
which plateau or even worsen past $N{\approx}21$ rather than continuing
to improve), consistent with the main paper's own finding (Section~6.8)
that the operator's error is dominated by optimization error rather than
discretization error -- it does not inherit FEM's guarantee of monotone
convergence under refinement, because refining $N$ only changes the
problem's discretization, not the fixed training budget the network was
optimized under.

\FloatBarrier
\section{Accuracy-threshold tables (1\%/2\%/5\%), all six cases}
\label{sec:s3}
The main paper's Table 21 reports only the $5\%$-threshold row, on each
case's own binding quantity ($H^1$ semi-norm in every case). The report's
own round-13 analysis tabulates, per quantity of interest and per case,
the coarsest GPU-native-FEM resolution reaching each of three thresholds
(1\%, 2\%, 5\%), that resolution's own measured cost, and the resulting
break-even against the operator's compile+TF32 cost at its $N{=}1401$
deployment resolution ($\approx394$~ms, case-independent -- the same
baseline used throughout the main paper's Section 6.3/6.6/6.7).
For some quantities the $1\%$ threshold is not reached by $N{=}49$, the
coarsest-to-finest range tested; those cells are marked n/a, exactly as
in the source.

\begin{table}[h]
\centering
\caption{B1$\times$Neo-Hookean: minimum GPU-native-FEM resolution $N$
  required to reach each QoI-accuracy threshold, that resolution's own
  real per-sample GPU-FEM cost, and the operator's break-even point
  against it.}
\label{tab:s3-b1nh}
\small
\begin{tabular}{lcccc}
\toprule
QoI & Threshold & Required $N$ & FEM (ms/sample) & Break-even (samples) \\
\midrule
Displacement $L^2$  & 1\% & 9   & 1{,}659.8 & 33{,}089 \\
Displacement $L^2$  & 2\% & 6   & 1{,}669.9 & 32{,}828 \\
Displacement $L^2$  & 5\% & 4   & 1{,}654.0 & 33{,}241 \\
$H^1$ semi-norm     & 1\% & n/a & n/a       & not reached by $N{=}49$ \\
$H^1$ semi-norm     & 2\% & 41  & 6{,}506.4 & 6{,}852 \\
$H^1$ semi-norm     & 5\% & 17  & 1{,}700.3 & 32{,}063 \\
Tangent energy      & 1\% & n/a & n/a       & not reached by $N{=}49$ \\
Tangent energy      & 2\% & 37  & 4{,}588.9 & 9{,}984 \\
Tangent energy      & 5\% & 13  & 1{,}672.0 & 32{,}773 \\
Reaction (B1 only)  & 1\% & 9   & 1{,}659.8 & 33{,}089 \\
Reaction (B1 only)  & 2\% & 9   & 1{,}659.8 & 33{,}089 \\
Reaction (B1 only)  & 5\% & 5   & 1{,}666.6 & 32{,}911 \\
Region-Cauchy avg.\ & 1\% & 21  & 1{,}736.9 & 31{,}189 \\
Region-Cauchy avg.\ & 2\% & 11  & 1{,}660.7 & 33{,}065 \\
Region-Cauchy avg.\ & 5\% & 3   & 1{,}636.0 & 33{,}723 \\
\bottomrule
\end{tabular}
\end{table}

\begin{table}[h]
\centering
\caption{B1$\times$Mooney-Rivlin: same protocol as Table~\ref{tab:s3-b1nh}.}
\label{tab:s3-b1mr}
\small
\begin{tabular}{lcccc}
\toprule
QoI & Threshold & Required $N$ & FEM (ms/sample) & Break-even (samples) \\
\midrule
Displacement $L^2$  & 1\% & 9   & 1{,}936.0 & 34{,}580 \\
Displacement $L^2$  & 2\% & 6   & 1{,}938.4 & 34{,}525 \\
Displacement $L^2$  & 5\% & 4   & 1{,}922.5 & 34{,}886 \\
$H^1$ semi-norm     & 1\% & n/a & n/a       & not reached by $N{=}49$ \\
$H^1$ semi-norm     & 2\% & 41  & 6{,}610.0 & 8{,}577 \\
$H^1$ semi-norm     & 5\% & 17  & 1{,}994.5 & 33{,}317 \\
Tangent energy      & 1\% & n/a & n/a       & not reached by $N{=}49$ \\
Tangent energy      & 2\% & 37  & 4{,}679.5 & 12{,}442 \\
Tangent energy      & 5\% & 13  & 1{,}969.6 & 33{,}843 \\
Reaction (B1 only)  & 1\% & 9   & 1{,}936.0 & 34{,}580 \\
Reaction (B1 only)  & 2\% & 6   & 1{,}938.4 & 34{,}525 \\
Reaction (B1 only)  & 5\% & 4   & 1{,}922.5 & 34{,}886 \\
Region-Cauchy avg.\ & 1\% & 25  & 2{,}099.0 & 31{,}274 \\
Region-Cauchy avg.\ & 2\% & 11  & 1{,}952.0 & 34{,}225 \\
Region-Cauchy avg.\ & 5\% & 9   & 1{,}936.0 & 34{,}580 \\
\bottomrule
\end{tabular}
\end{table}

\begin{table}[h]
\centering
\caption{B1$\times$Arruda-Boyce: same protocol as Table~\ref{tab:s3-b1nh}.}
\label{tab:s3-b1ab}
\small
\begin{tabular}{lcccc}
\toprule
QoI & Threshold & Required $N$ & FEM (ms/sample) & Break-even (samples) \\
\midrule
Displacement $L^2$  & 1\% & 9   & 2{,}637.9 & 17{,}377 \\
Displacement $L^2$  & 2\% & 5   & 2{,}585.9 & 17{,}789 \\
Displacement $L^2$  & 5\% & 4   & 2{,}596.5 & 17{,}704 \\
$H^1$ semi-norm     & 1\% & n/a & n/a       & not reached by $N{=}49$ \\
$H^1$ semi-norm     & 2\% & 33  & 3{,}522.2 & 12{,}464 \\
$H^1$ semi-norm     & 5\% & 13  & 2{,}677.1 & 17{,}078 \\
Tangent energy      & 1\% & n/a & n/a       & not reached by $N{=}49$ \\
Tangent energy      & 2\% & 29  & 2{,}822.4 & 16{,}056 \\
Tangent energy      & 5\% & 11  & 2{,}636.1 & 17{,}391 \\
Reaction (B1 only)  & 1\% & 9   & 2{,}637.9 & 17{,}377 \\
Reaction (B1 only)  & 2\% & 9   & 2{,}637.9 & 17{,}377 \\
Reaction (B1 only)  & 5\% & 5   & 2{,}585.9 & 17{,}789 \\
Region-Cauchy avg.\ & 1\% & 4   & 2{,}596.5 & 17{,}704 \\
Region-Cauchy avg.\ & 2\% & 3   & 2{,}590.9 & 17{,}748 \\
Region-Cauchy avg.\ & 5\% & 3   & 2{,}590.9 & 17{,}748 \\
\bottomrule
\end{tabular}
\end{table}

\begin{table}[h]
\centering
\caption{B2$\times$Neo-Hookean: same protocol as Table~\ref{tab:s3-b1nh}
  (B2 has no single reaction-resultant metric, main paper Table 20).}
\label{tab:s3-b2nh}
\small
\begin{tabular}{lcccc}
\toprule
QoI & Threshold & Required $N$ & FEM (ms/sample) & Break-even (samples) \\
\midrule
Displacement $L^2$  & 1\% & 13  & 1{,}710.3 & 30{,}053 \\
Displacement $L^2$  & 2\% & 9   & 1{,}669.4 & 31{,}016 \\
Displacement $L^2$  & 5\% & 6   & 1{,}652.1 & 31{,}442 \\
$H^1$ semi-norm     & 1\% & n/a & n/a       & not reached by $N{=}49$ \\
$H^1$ semi-norm     & 2\% & 41  & 6{,}524.0 & 6{,}454 \\
$H^1$ semi-norm     & 5\% & 17  & 1{,}728.9 & 29{,}633 \\
Tangent energy      & 1\% & n/a & n/a       & not reached by $N{=}49$ \\
Tangent energy      & 2\% & 33  & 3{,}199.5 & 14{,}101 \\
Tangent energy      & 5\% & 17  & 1{,}728.9 & 29{,}633 \\
Region-Cauchy avg.\ & 1\% & 13  & 1{,}710.3 & 30{,}053 \\
Region-Cauchy avg.\ & 2\% & 9   & 1{,}669.4 & 31{,}016 \\
Region-Cauchy avg.\ & 5\% & 6   & 1{,}652.1 & 31{,}442 \\
\bottomrule
\end{tabular}
\end{table}

\begin{table}[h]
\centering
\caption{B2$\times$Mooney-Rivlin: same protocol as Table~\ref{tab:s3-b1nh}.}
\label{tab:s3-b2mr}
\small
\begin{tabular}{lcccc}
\toprule
QoI & Threshold & Required $N$ & FEM (ms/sample) & Break-even (samples) \\
\midrule
Displacement $L^2$  & 1\% & 11  & 1{,}985.7 & 3{,}982 \\
Displacement $L^2$  & 2\% & 9   & 1{,}985.0 & 3{,}983 \\
Displacement $L^2$  & 5\% & 6   & 1{,}961.3 & 4{,}043 \\
$H^1$ semi-norm     & 1\% & n/a & n/a       & not reached by $N{=}49$ \\
$H^1$ semi-norm     & 2\% & 41  & 6{,}630.8 & 1{,}016 \\
$H^1$ semi-norm     & 5\% & 17  & 2{,}040.4 & 3{,}849 \\
Tangent energy      & 1\% & n/a & n/a       & not reached by $N{=}49$ \\
Tangent energy      & 2\% & 33  & 3{,}243.9 & 2{,}223 \\
Tangent energy      & 5\% & 17  & 2{,}040.4 & 3{,}849 \\
Region-Cauchy avg.\ & 1\% & 13  & 1{,}996.0 & 3{,}956 \\
Region-Cauchy avg.\ & 2\% & 9   & 1{,}985.0 & 3{,}983 \\
Region-Cauchy avg.\ & 5\% & 6   & 1{,}961.3 & 4{,}043 \\
\bottomrule
\end{tabular}
\end{table}

\begin{table}[h]
\centering
\caption{B2$\times$Arruda-Boyce: same protocol as Table~\ref{tab:s3-b1nh}.}
\label{tab:s3-b2ab}
\small
\begin{tabular}{lcccc}
\toprule
QoI & Threshold & Required $N$ & FEM (ms/sample) & Break-even (samples) \\
\midrule
Displacement $L^2$  & 1\% & 13  & 2{,}652.3 & 3{,}572 \\
Displacement $L^2$  & 2\% & 9   & 2{,}648.4 & 3{,}578 \\
Displacement $L^2$  & 5\% & 6   & 2{,}593.3 & 3{,}668 \\
$H^1$ semi-norm     & 1\% & n/a & n/a       & not reached by $N{=}49$ \\
$H^1$ semi-norm     & 2\% & 41  & 7{,}091.9 & 1{,}204 \\
$H^1$ semi-norm     & 5\% & 17  & 2{,}666.0 & 3{,}551 \\
Tangent energy      & 1\% & n/a & n/a       & not reached by $N{=}49$ \\
Tangent energy      & 2\% & 33  & 3{,}545.3 & 2{,}560 \\
Tangent energy      & 5\% & 17  & 2{,}666.0 & 3{,}551 \\
Region-Cauchy avg.\ & 1\% & 13  & 2{,}652.3 & 3{,}572 \\
Region-Cauchy avg.\ & 2\% & 9   & 2{,}648.4 & 3{,}578 \\
Region-Cauchy avg.\ & 5\% & 6   & 2{,}593.3 & 3{,}668 \\
\bottomrule
\end{tabular}
\end{table}

A pattern worth stating plainly, per the report's own summary: across
all six cases and every threshold tested, the required FEM resolution
never exceeds $N{=}41$, and its own per-sample cost there (up to
$\approx7.1$~s) never comes close to the operator's $\approx394$~ms --
FEM never becomes cheaper than the operator at any accuracy level tested,
for any case or quantity.

\FloatBarrier
\section{Multi-resolution training-cost breakdown}
\label{sec:s4}
Main paper Section 6.5 (zero-shot resolution invariance, B1$\times$Neo-Hookean)
reports that a direct-$N{=}1401$-trained checkpoint is cheaper to train
than the multi-resolution checkpoint but substantially less accurate at
unseen resolutions. Table~\ref{tab:s4} reports the underlying per-run
training-cost breakdown for all seven completed training runs this
comparison and the main results draw on, read from each run's own
training-metrics log.

\begin{table}[h]
\centering
\caption{Training-cost summary, all seven completed training runs:
  resolutions used, samples per resolution, final epoch/optimizer steps,
  wall-clock, and cost per sample/optimizer step. B2$\times$Arruda-Boyce
  has no multi-resolution retrain; its original fixed-resolution
  ($N{=}21,33$) run is reported instead.}
\label{tab:s4}
\footnotesize
\begin{tabular}{lcccccc}
\toprule
Run & Resolutions & Samples/res & Final epoch & Opt.\ steps & Wall-clock & Cost/sample \\
\midrule
B1$\times$NH (multi-res)       & 21,33,101,201 & 400 & 1{,}075 & 215{,}000 & 11.63~h & 26.18~s \\
B1$\times$NH (direct-$N{=}1401$) & 1401        & 100 & 36      & 3{,}600   & 8.19~h  & 294.85~s \\
B1$\times$MR (multi-res)       & 21,33,101,201 & 400 & 1{,}350 & 270{,}000 & 14.81~h & 33.31~s \\
B1$\times$AB (multi-res)       & 21,33,101,201 & 400 & 975     & 195{,}000 & 10.83~h & 24.36~s \\
B2$\times$NH (multi-res)       & 21,33,101,201 & 400 & 1{,}875 & 375{,}000 & 10.99~h & 24.73~s \\
B2$\times$MR (multi-res)       & 21,33,101,201 & 400 & 300     & 60{,}000  & 1.76~h  & 3.95~s \\
B2$\times$AB (original, no retrain) & 21,33     & 400 & 2{,}200 & 220{,}000 & 2.24~h  & 10.06~s \\
\bottomrule
\end{tabular}
\end{table}

The direct-$N{=}1401$ run's $294.85$~s/sample is an order of magnitude
above every multi-resolution run's own per-sample cost ($3.95$--$33.31$~s),
not just the $21\%$ difference the main paper's own B1$\times$NH-specific
comparison quotes -- that $21\%$ figure compares total training
wall-clock for that one case (direct: $8.19$~h vs.\ multi-res: $11.63$~h),
not per-sample cost, since the direct run uses far fewer total samples.

\FloatBarrier
\section{Throughput and memory-ceiling comparison}
\label{sec:s5}
Main paper Section 6.3 states the operator's peak throughput exceeds the
native finite-element solver's by roughly $1{,}470\times$ at matched GPU
memory, with Figure 25 showing the trend. Table~\ref{tab:s5} reports the
full underlying measurement (B1$\times$Neo-Hookean, $N{=}21$, same GPU,
corrected checkpoint).

\begin{table}[h]
\centering
\caption{Max feasible batch size and throughput, operator vs.\ native
  finite-element solver, same GPU, $N{=}21$.}
\label{tab:s5}
\small
\begin{tabular}{lccc}
\toprule
Scenario & Max batch & Peak GPU mem. & Throughput (samples/s) \\
\midrule
Operator, own ceiling      & 8{,}192 & 47.6~GB & 4{,}203.79 \\
FEM, own ceiling           & 256     & 38.7~GB & 2.86 \\
FEM, capped at op.\ mem.\  & 256     & 38.7~GB & 2.86 \\
Operator, capped at FEM mem.\ & 4{,}096 & 23.8~GB & 4{,}207.75 \\
\bottomrule
\end{tabular}
\end{table}

The finite-element solver already saturates below either method's memory
ceiling (its own maximum batch size, 256, uses only $38.7$~GB of the
$47.6$--$80$~GB available), while the operator continues scaling batch
size well past that point with no loss of throughput -- the throughput
gap is $\approx1{,}470\times$ regardless of which memory budget is used
to cap the comparison.

\FloatBarrier
\section{Inference-speed optimization breakdown}
\label{sec:s6}
Main paper Section 7.4 and Figure 22 use the operator's compile+TF32
inference cost at $N{=}1401$ ($\approx394$~ms) as the baseline for every
B3 accuracy-matched break-even figure reported. Table~\ref{tab:s6} reports
the full optimization-attempt breakdown this number comes from, including
each variant's numerical output difference against the unoptimized
eager-FP32 baseline (confirming the optimizations are numerically benign,
not merely fast).

\begin{table}[h]
\centering
\caption{Inference-speed optimization attempts, $N{=}1401$, batch size 1,
  corrected checkpoint.}
\label{tab:s6}
\small
\begin{tabular}{lccc}
\toprule
Variant & ms/sample & Speed-up vs.\ eager & Output diff.\ vs.\ eager \\
\midrule
Eager (FP32)      & 2{,}292.1 & $1.00\times$ & -- \\
\texttt{torch.compile}        & 2{,}145.1 & $1.07\times$ & $6.0\times10^{-7}$ \\
Eager + TF32      & 491.2     & $4.67\times$ & $3.2\times10^{-3}$ \\
Compile + TF32    & 394.0     & $5.82\times$ & $3.2\times10^{-3}$ \\
\bottomrule
\end{tabular}
\end{table}

\FloatBarrier
\section{Progressive out-of-distribution shift}
\label{sec:s7}
Main paper Section 6.4 reports out-of-distribution degradation only at
the full combined-shift endpoint used throughout the study
($2.5\sigma$ material / $2$--$2.5\sigma$ loading, Table 16/17). The report's
own progressive sweep, isolating material and loading shifts separately
across six intermediate shift magnitudes on B1$\times$Neo-Hookean, is
reported in Table~\ref{tab:s7a}, and a companion raw-input-vs-normalized-input
ablation -- testing whether standardizing B1's own four input channels
(a separate normalization from the B3-specific input-normalization fix
discussed in the main paper's Section~7.1, unrelated to it) changes the
OOD sensitivity picture -- in Table~\ref{tab:s7b}.

\begin{table}[h]
\centering
\caption{Progressive out-of-distribution shift, B1$\times$Neo-Hookean,
  one factor at a time. Each entry is the mean relative $L^2$ error over
  ten held-out samples, with its ratio ($\times$) to the in-distribution
  baseline ($k{=}0$: $0.0867$).}
\label{tab:s7a}
\small
\begin{tabular}{lcccccc}
\toprule
Shift $k$ ($\sigma$) & Loading & $\times$ & Material & $\times$ & Both & $\times$ \\
\midrule
0.5 & 0.0869 & $1.00\times$ & 0.1409 & $1.63\times$ & 0.1427 & $1.65\times$ \\
1.0 & 0.0888 & $1.02\times$ & 0.2073 & $2.39\times$ & 0.2066 & $2.38\times$ \\
1.5 & 0.0897 & $1.03\times$ & 0.2763 & $3.19\times$ & 0.2682 & $3.09\times$ \\
2.0 & 0.0871 & $1.00\times$ & 0.3486 & $4.02\times$ & 0.3294 & $3.80\times$ \\
2.5 & 0.0862 & $0.99\times$ & 0.4265 & $4.92\times$ & 0.3957 & $4.56\times$ \\
3.0 & 0.0927 & $1.07\times$ & 0.5112 & $5.90\times$ & 0.4601 & $5.31\times$ \\
\bottomrule
\end{tabular}
\end{table}

\begin{table}[h]
\centering
\caption{Raw-input vs.\ standardized-input OOD sensitivity, B1$\times$Neo-Hookean
  (same progressive sweep as Table~\ref{tab:s7a}, two otherwise-identical
  checkpoints). Each model's own ``$\times$'' column is normalized against
  that same model's own in-distribution baseline, not the other model's.}
\label{tab:s7b}
\small
\begin{tabular}{llcccc}
\toprule
Factor & Shift $k$ & Raw error & Normalized error & Raw $\times$ & Normalized $\times$ \\
\midrule
Material & 0.5 & 0.1409 & 0.1700 & $1.63\times$ & $1.84\times$ \\
Material & 1.0 & 0.2073 & 0.2577 & $2.39\times$ & $2.79\times$ \\
Material & 1.5 & 0.2763 & 0.3385 & $3.19\times$ & $3.66\times$ \\
Material & 2.0 & 0.3486 & 0.3973 & $4.02\times$ & $4.30\times$ \\
Material & 2.5 & 0.4265 & 0.4030 & $4.92\times$ & $4.36\times$ \\
Material & 3.0 & 0.5112 & 0.3565 & $5.90\times$ & $3.85\times$ \\
\midrule
Loading  & 0.5 & 0.0869 & 0.0992 & $1.00\times$ & $1.07\times$ \\
Loading  & 1.0 & 0.0888 & 0.1025 & $1.02\times$ & $1.11\times$ \\
Loading  & 1.5 & 0.0897 & 0.1043 & $1.03\times$ & $1.13\times$ \\
Loading  & 2.0 & 0.0871 & 0.1040 & $1.00\times$ & $1.12\times$ \\
Loading  & 2.5 & 0.0862 & 0.1015 & $0.99\times$ & $1.10\times$ \\
Loading  & 3.0 & 0.0927 & 0.1274 & $1.07\times$ & $1.38\times$ \\
\midrule
Both     & 0.5 & 0.1427 & 0.1747 & $1.65\times$ & $1.89\times$ \\
Both     & 1.0 & 0.2066 & 0.2459 & $2.38\times$ & $2.66\times$ \\
Both     & 1.5 & 0.2682 & 0.2979 & $3.09\times$ & $3.22\times$ \\
Both     & 2.0 & 0.3294 & 0.3293 & $3.80\times$ & $3.56\times$ \\
Both     & 2.5 & 0.3957 & 0.3540 & $4.56\times$ & $3.83\times$ \\
Both     & 3.0 & 0.4601 & 0.4064 & $5.31\times$ & $4.39\times$ \\
\bottomrule
\end{tabular}
\end{table}

The material shift dominates at every magnitude tested, consistent with
the combined-endpoint isolation already reported in the main paper's
Table 17; the loading-only column stays close to flat ($0.99$--$1.07\times$)
across the entire range, while the material-only column grows
monotonically and far faster. Input standardization does not change this
qualitative picture (material still dominates loading in the normalized
model), though it does shift the absolute sensitivity at the largest
material shifts tested.

\FloatBarrier
\section{Full manufactured-solution operator-vs-Q4/Q9 tables}
\label{sec:s8}
Main paper Section 6.8 and its Table~25 report a compact two-part summary
(single-member absolute errors at three meshes, and 16-member family-mean
operator/Q4 ratios in all four norms) supporting the claim that the
operator's own convergence rate is negative ($-0.59$ in $L^2$) while Q4's
is positive ($+1.99$). The full underlying per-norm, per-mesh tables are
reported here for completeness.

\begin{table}[h]
\centering
\caption{Operator vs.\ Q4 vs.\ Q9, single manufactured member, $N{=}17$:
  the three-way comparison every other table in this section is anchored
  to (same routine, same manufactured solution).}
\label{tab:s8-table24}
\small
\begin{tabular}{lccccc}
\toprule
Method & DOF & $L^2$ & $H^1$ semi-norm & Stress & Energy \\
\midrule
Q4 (same mesh)              & 578   & $3.403\times10^{-3}$ & $5.666\times10^{-2}$ & $5.724\times10^{-2}$ & $3.191\times10^{-3}$ \\
Q9 (same $N$)                & 2{,}178 & $5.163\times10^{-5}$ & $1.438\times10^{-3}$ & $1.495\times10^{-3}$ & $2.088\times10^{-6}$ \\
Physics-informed operator   & 578   & $8.238\times10^{-3}$ & $5.831\times10^{-2}$ & $5.886\times10^{-2}$ & $9.914\times10^{-3}$ \\
\bottomrule
\end{tabular}
\end{table}

\begin{table}[h]
\centering
\caption{Operator/Q4 error ratio, single manufactured member, all four
  norms, three meshes (values above $1\times$ mean the operator is
  further from the manufactured solution than Q4 at the same mesh).}
\label{tab:s8-table24b}
\small
\begin{tabular}{lcccc}
\toprule
$N$ & $L^2$ & $H^1$ semi-norm & Stress & Energy \\
\midrule
9  & $0.37\times$  & $1.01\times$ & $1.01\times$ & $0.89\times$ \\
17 & $2.42\times$  & $1.03\times$ & $1.03\times$ & $3.11\times$ \\
33 & $13.33\times$ & $1.36\times$ & $1.33\times$ & $6.62\times$ \\
\bottomrule
\end{tabular}
\end{table}

\begin{table}[h]
\centering
\caption{Q4, Q9 and the operator over the operator's own 16-member test
  family, relative $L^2$ against the manufactured solution (family
  means). The rightmost two columns compare the family-mean ratio
  against the single-member ratio of Table~\ref{tab:s8-table24b}.}
\label{tab:s8-table24c}
\footnotesize
\begin{tabular}{lccccc}
\toprule
$N$ & Q4 mean & Q9 mean & Operator mean & \multicolumn{2}{c}{Operator/Q4} \\
 & & & & family & single member \\
\midrule
9  & $1.3515\times10^{-2}$ & $4.1550\times10^{-4}$ & $8.3792\times10^{-3}$ & $0.62\times$  & $0.37\times$ \\
17 & $3.4049\times10^{-3}$ & $5.1628\times10^{-5}$ & $8.8261\times10^{-3}$ & $2.59\times$  & $2.42\times$ \\
33 & $8.5298\times10^{-4}$ & $6.4423\times10^{-6}$ & $1.2358\times10^{-2}$ & $14.49\times$ & $13.33\times$ \\
\bottomrule
\end{tabular}
\end{table}

\begin{table}[h]
\centering
\caption{Coefficient of variation (std/mean) across the 16-member family,
  operator against Q4, all four norms. Q4's own spread is negligible at
  every mesh (member-independent); the operator's is roughly two orders
  of magnitude larger in $L^2$ and energy, already at the coarsest mesh.}
\label{tab:s8-table24e}
\small
\begin{tabular}{lcccc}
\toprule
$N$ & $L^2$ & $H^1$ semi-norm & Stress & Energy \\
\midrule
9  & 0.384 (Q4 0.002) & 0.003 (Q4 0.000) & 0.006 (Q4 0.007) & 0.549 (Q4 0.003) \\
17 & 0.258 (Q4 0.003) & 0.016 (Q4 0.000) & 0.015 (Q4 0.007) & 0.370 (Q4 0.003) \\
33 & 0.425 (Q4 0.003) & 0.134 (Q4 0.000) & 0.120 (Q4 0.007) & 0.736 (Q4 0.003) \\
\bottomrule
\end{tabular}
\end{table}

\FloatBarrier
\section{Q4-vs-Q9 fine-reference agreement (PASS/FAIL)}
\label{sec:s9}
Main paper Section 4 states in prose that the Q4 and Q9 fine references
``agree to within $2\times10^{-6}$ in $L^2$ (though not fully to the
prescribed $10^{-5}$ criterion in $H^1$/energy).'' Table~\ref{tab:s9}
reports the direct PASS/FAIL comparison against the advisor's own stated
criterion this sentence summarizes.

\begin{table}[h]
\centering
\caption{Direct Q4-vs-Q9 fine-solution comparison (B1$\times$Neo-Hookean,
  $N{=}2236$, $\approx10$M/40M DOF) against the advisor's explicit
  $<10^{-5}$ criterion.}
\label{tab:s9}
\small
\begin{tabular}{lcccc}
\toprule
Norm & Rel.\ diff.\ (Q4-domain) & Rel.\ diff.\ (Q9-domain) & Target & Verdict \\
\midrule
$L^2$            & $2.06\times10^{-6}$ & $2.06\times10^{-6}$ & $<10^{-5}$ & PASS \\
$H^1$ semi-norm  & $1.11\times10^{-3}$ & $1.15\times10^{-3}$ & $<10^{-5}$ & FAIL \\
Tangent energy   & $9.61\times10^{-4}$ & $3.10\times10^{-4}$ & $<10^{-5}$ & FAIL \\
\bottomrule
\end{tabular}
\end{table}

The $H^1$ and energy ``FAIL'' verdicts are read here as a precision
statement about the two fine references, not a correctness concern for
any result built on them: both are several orders of magnitude tighter
than any accuracy figure reported elsewhere in the main paper, and the
$L^2$ displacement field itself -- the quantity every downstream
comparison is actually scored against -- passes the criterion directly.

\FloatBarrier
\section{GPU training-memory footprint}
\label{sec:s10}
Main paper Figure 7 plots training memory footprint across all six
cases without a supporting table. Table~\ref{tab:s10} reports the three
measured levels directly: PyTorch-allocated, PyTorch-reserved, and
device-level (\texttt{nvidia-smi}-equivalent) peak memory, each from a
short, dedicated 3-epoch re-run per case (same architecture, same batch
size 8) performed specifically to capture this third instrumentation
channel, which predates the original six training runs.

\begin{table}[h]
\centering
\caption{Peak GPU memory during training, all six cases, batch size 8.}
\label{tab:s10}
\small
\begin{tabular}{lccc}
\toprule
Case & Allocated (MB) & Reserved (MB) & Device-level (MB) \\
\midrule
B1 $\times$ NH & 425.05 & 679.48 & 1{,}226.77 \\
B1 $\times$ MR & 425.05 & 677.38 & 1{,}224.67 \\
B1 $\times$ AB & 425.05 & 679.48 & 1{,}228.87 \\
B2 $\times$ NH & 425.05 & 679.48 & 1{,}226.77 \\
B2 $\times$ MR & 425.05 & 677.38 & 1{,}224.67 \\
B2 $\times$ AB & 425.05 & 679.48 & 1{,}228.87 \\
\bottomrule
\end{tabular}
\end{table}

Peak allocated memory is essentially identical across all six cases
($425.05$~MB, consistent with all six sharing the same architecture and
batch size); device-level peak ($\approx1.225$~GB) is only about
$1.4\%$ of an 80~GB A100's available memory -- training memory is not a
binding constraint for any of the six 2D benchmark cases.

\end{document}
